
\documentclass[11pt]{article}
\usepackage[letterpaper,margin=0.72in]{geometry}
\usepackage[T1]{fontenc}
\usepackage[utf8]{inputenc}
\usepackage{lmodern}
\usepackage{array}
\usepackage{booktabs}
\usepackage{longtable}
\usepackage{tabularx}
\usepackage{enumitem}
\usepackage{xcolor}
\usepackage{hyperref}
\usepackage{fancyhdr}
\usepackage{lastpage}
\usepackage{pdflscape}
\usepackage{listings}
\hypersetup{colorlinks=true,linkcolor=black,urlcolor=blue,citecolor=black}
\definecolor{HeaderBlue}{HTML}{1F4E79}
\definecolor{LightBlue}{HTML}{EAF3F8}
\definecolor{LightGray}{HTML}{F5F7FA}
\definecolor{RuleGray}{HTML}{B7C3D0}
\definecolor{RiskAmber}{HTML}{FFF4D6}
\newcolumntype{Y}{>{\raggedright\arraybackslash}X}
\newcolumntype{L}[1]{>{\raggedright\arraybackslash}p{#1}}
\setlength{\parindent}{0pt}
\setlength{\parskip}{0.45em}
\setlist[itemize]{leftmargin=1.2em,itemsep=0.18em,topsep=0.2em}
\setlist[enumerate]{leftmargin=1.4em,itemsep=0.18em,topsep=0.2em}
\pagestyle{fancy}
\setlength{\headheight}{14pt}
\fancyhf{}
\lhead{Gauntlt DevSecOps Security Testing}
\rhead{\thepage\ of \pageref{LastPage}}
\renewcommand{\headrulewidth}{0.4pt}
\sloppy
\lstset{basicstyle=\ttfamily\small,breaklines=true,columns=fullflexible,keepspaces=true,frame=single,framerule=0.3pt,rulecolor=\color{RuleGray},backgroundcolor=\color{LightGray}}
\newcommand{\TitleBlock}[2]{%
  \begin{center}
  {\LARGE\bfseries\color{HeaderBlue} #1}\\[0.35em]
  {\large #2}\\[0.4em]
  {\small Prepared for DevSecOps implementation planning | Updated May 6, 2026}
  \end{center}
  \vspace{0.4em}\hrule\vspace{0.8em}
}
\newsavebox{\gauntltbox}
\newenvironment{Callout}[1]{%
  \par\vspace{0.3em}\noindent\begin{lrbox}{\gauntltbox}\begin{minipage}{0.96\linewidth}\textbf{#1}\par\smallskip}
  {\end{minipage}\end{lrbox}\fcolorbox{RuleGray}{LightBlue}{\usebox{\gauntltbox}}\par\vspace{0.4em}}
\newenvironment{RiskBox}[1]{%
  \par\vspace{0.3em}\noindent\begin{lrbox}{\gauntltbox}\begin{minipage}{0.96\linewidth}\textbf{#1}\par\smallskip}
  {\end{minipage}\end{lrbox}\fcolorbox{RuleGray}{RiskAmber}{\usebox{\gauntltbox}}\par\vspace{0.4em}}
\newcommand{\checkbox}{\raisebox{0.2ex}{\fbox{\rule{0pt}{0.9ex}\rule{0.9ex}{0pt}}}}
\begin{document}

\TitleBlock{Gauntlt Implementation Checklist}{Minimum checklist for adopting Gauntlt in CI/CD}

\begin{Callout}{How to Use This Checklist}
Use this checklist as a readiness worksheet before adding Gauntlt attacks to a delivery pipeline. Each item should have an owner, evidence location, and decision status.
\end{Callout}

\section{1. Governance and Authorization}
\begin{itemize}
  \item \checkbox\ Document the authorized test scope: applications, URLs, environments, ports, and excluded systems.
  \item \checkbox\ Confirm written permission for any dynamic tests, even when tests are non-destructive.
  \item \checkbox\ Define which attacks are allowed in CI, which require security approval, and which are prohibited.
  \item \checkbox\ Establish a stop condition for noisy, destructive, or high-impact tests.
  \item \checkbox\ Map each attack file to a control objective, user story, or security requirement.
\end{itemize}

\section{2. Environment Readiness}
\begin{itemize}
  \item \checkbox\ Identify a staging, ephemeral, or security test environment that mirrors production-relevant controls.
  \item \checkbox\ Confirm test accounts, seeded data, and rate limits are safe for repeatable automated testing.
  \item \checkbox\ Define allowlists for CI runners, scanners, and target services.
  \item \checkbox\ Confirm logging and monitoring can distinguish approved security tests from suspicious traffic.
  \item \checkbox\ Create rollback or cleanup steps for any test that changes target state.
\end{itemize}

\section{3. Toolchain Dependencies}
\begin{itemize}
  \item \checkbox\ Install Ruby and Gauntlt using a reproducible runner image or build step.
  \item \checkbox\ Install required adapters/tools such as \texttt{curl}, \texttt{nmap}, \texttt{sqlmap}, \texttt{dirb}, \texttt{SSLyze}, or project-specific scripts.
  \item \checkbox\ Pin versions for Gauntlt and external tools when feasible.
  \item \checkbox\ Run a dependency validation attack before running target-specific attacks.
  \item \checkbox\ Store tool installation evidence in the pipeline logs or build artifact metadata.
\end{itemize}

\section{4. Attack File Design}
\begin{itemize}
  \item \checkbox\ Keep one attack file focused on one security objective or closely related objective group.
  \item \checkbox\ Use readable feature and scenario names.
  \item \checkbox\ Separate preconditions, attack command, and expected result clearly.
  \item \checkbox\ Avoid brittle assertions that depend on unstable timestamps, random IDs, or unrelated response text.
  \item \checkbox\ Include negative checks for unacceptable findings where supported by the adapter.
\end{itemize}

\section{5. CI/CD Integration}
\begin{itemize}
  \item \checkbox\ Run Gauntlt after deployment to the test environment, not only during source build.
  \item \checkbox\ Fail the job when required attacks fail.
  \item \checkbox\ Publish machine-readable and human-readable results where possible.
  \item \checkbox\ Archive attack files, command output, target version, commit SHA, and environment ID as evidence.
  \item \checkbox\ Route failures to the owning development team with reproduction steps.
\end{itemize}

\section{6. Quality Gates and Exception Handling}
\begin{itemize}
  \item \checkbox\ Define severity thresholds for blocking release.
  \item \checkbox\ Define who can approve a temporary exception.
  \item \checkbox\ Require expiration dates for exceptions.
  \item \checkbox\ Require compensating controls when a release proceeds despite a failed security check.
  \item \checkbox\ Re-run failed attacks after remediation before closing the issue.
\end{itemize}

\section{7. Evidence Checklist}
\begin{longtable}{@{}L{0.30\linewidth}L{0.24\linewidth}L{0.36\linewidth}@{}}
\toprule
\textbf{Evidence Item} & \textbf{Required?} & \textbf{Notes} \\
\midrule
Attack file path & Yes & Store in version control. \\
Pipeline run URL & Yes & Link to the exact CI execution. \\
Target environment & Yes & Include URL, environment name, and deployment ID. \\
Command output & Yes & Preserve stdout, stderr, and exit code. \\
Tool versions & Recommended & Helps explain drift and reproduce failures. \\
Exception record & Conditional & Required when a failed gate is bypassed. \\
Remediation ticket & Conditional & Required when failures create follow-up work. \\
\bottomrule
\end{longtable}


\section*{References}
\begin{itemize}
  \item Gauntlt official site: \url{http://gauntlt.org/}
  \item Gauntlt GitHub repository: \url{https://github.com/gauntlt/gauntlt}
  \item Gauntlt Starter Kit: \url{https://github.com/gauntlt/gauntlt-starter-kit}
  \item Gauntlt Demo: \url{https://github.com/gauntlt/gauntlt-demo}
\end{itemize}

\end{document}



